%%%%%%%%%%%%%%%%%%%%%%%%%%%%%%%%%%%%%%%%%%%%%%%%%%%%%%%%%%%%%%%%%%%
% Standalone general university LaTeX template -- homework sheet
% One-file version: no \input files needed.
%%%%%%%%%%%%%%%%%%%%%%%%%%%%%%%%%%%%%%%%%%%%%%%%%%%%%%%%%%%%%%%%%%%

\documentclass[11pt,a4paper]{scrartcl}

%%%%%%%%%%%%%%%%%%%%%%%%%%%%%%%%%%%%%%%%%%%%%%%%%%%%%%%%%%%%%%%%%%%
% Embedded file: includes/university_metadata.tex
%%%%%%%%%%%%%%%%%%%%%%%%%%%%%%%%%%%%%%%%%%%%%%%%%%%%%%%%%%%%%%%%%%%
%%%%%%%%%%%%%%%%%%%%%%%%%%%%%%%%%%%%%%%%%%%%%%%%%%%%%%%%%%%%%%%%%%%
% General university LaTeX template -- metadata and title pages
%%%%%%%%%%%%%%%%%%%%%%%%%%%%%%%%%%%%%%%%%%%%%%%%%%%%%%%%%%%%%%%%%%%

% Change these in the main .tex file before \begin{document}.
\newcommand{\CourseTitle}{COURSE NAME}
\newcommand{\CourseShortTitle}{COURSE SHORT NAME}
\newcommand{\DocumentKind}{DOCUMENT KIND}
\newcommand{\DocumentTitle}{DOCUMENT TITLE}
\newcommand{\DocumentSubtitle}{}
\newcommand{\AuthorName}{YOUR NAME}
\newcommand{\InstitutionName}{INSTITUTION NAME}
\newcommand{\SubmissionDate}{\today}

\newcommand{\makeuniversitytitle}{%
    \begin{titlepage}
        \centering
        \vspace*{2cm}
        {\Large \CourseTitle\par}
        \vspace{1.2cm}
        {\huge\bfseries \DocumentTitle\par}
        \ifx\DocumentSubtitle\empty\else
            \vspace{0.4cm}
            {\Large \DocumentSubtitle\par}
        \fi
        \vspace{1cm}
        {\Large \DocumentKind\par}
        \vfill
        {\large \AuthorName\par}
        \vspace{0.2cm}
        {\large \InstitutionName\par}
        \vspace{0.8cm}
        {\large \SubmissionDate\par}
    \end{titlepage}
    \pagenumbering{arabic}
}

\newcommand{\makechaptertitle}{%
    \begin{center}
        {\Large \CourseTitle\par}
        \vspace{0.5cm}
        {\huge\bfseries \DocumentTitle\par}
        \ifx\DocumentSubtitle\empty\else
            \vspace{0.25cm}
            {\Large \DocumentSubtitle\par}
        \fi
        \vspace{0.5cm}
        {\large \AuthorName \hfill \SubmissionDate\par}
    \end{center}
    \vspace{0.5cm}
}

% Backward-compatible alias for older documents created with the Physik IV template.
\newcommand{\makephysiktitle}{\makeuniversitytitle}

%%%%%%%%%%%%%%%%%%%%%%%%%%%%%%%%%%%%%%%%%%%%%%%%%%%%%%%%%%%%%%%%%%%
% Document metadata -- change these placeholders for each sheet
%%%%%%%%%%%%%%%%%%%%%%%%%%%%%%%%%%%%%%%%%%%%%%%%%%%%%%%%%%%%%%%%%%%
\renewcommand{\CourseTitle}{Physik IV: Kerne und Teilchen}
\renewcommand{\CourseShortTitle}{Physik IV}
\renewcommand{\DocumentKind}{Hausaufgabe, Tafelvortrag und Lernfassung}
\renewcommand{\DocumentTitle}{Aufgabenblatt 9}
\renewcommand{\DocumentSubtitle}{Relativistische Kinematik, Schwerpunktsystem und invariante Masse}
\renewcommand{\AuthorName}{Tobias Laser}
\renewcommand{\InstitutionName}{Universit\"at Innsbruck}
\renewcommand{\SubmissionDate}{26. Mai 2026}

%%%%%%%%%%%%%%%%%%%%%%%%%%%%%%%%%%%%%%%%%%%%%%%%%%%%%%%%%%%%%%%%%%%
% Embedded file: includes/university_packages.tex
%%%%%%%%%%%%%%%%%%%%%%%%%%%%%%%%%%%%%%%%%%%%%%%%%%%%%%%%%%%%%%%%%%%
%%%%%%%%%%%%%%%%%%%%%%%%%%%%%%%%%%%%%%%%%%%%%%%%%%%%%%%%%%%%%%%%%%%
% General university LaTeX template -- packages and global setup
%%%%%%%%%%%%%%%%%%%%%%%%%%%%%%%%%%%%%%%%%%%%%%%%%%%%%%%%%%%%%%%%%%%

% Page layout and spacing
\usepackage[margin=2.5cm]{geometry}
\usepackage[onehalfspacing]{setspace}
\usepackage{scrlayer-scrpage}

% Font and language
\usepackage[T1]{fontenc}
\usepackage[utf8]{inputenc}
\usepackage[ngerman]{babel}
\usepackage{lmodern}
\usepackage{microtype}
\usepackage{csquotes}

% Mathematics
\usepackage{amsmath,amsthm,amssymb,mathtools}
\usepackage{bm}
\usepackage{icomma}
\usepackage{cancel}

% Graphics, plots, diagrams
\usepackage{graphicx}
\usepackage{tikz}
\usetikzlibrary{arrows.meta,calc,decorations.pathmorphing,positioning,patterns,angles,quotes}
\usepackage{pgfplots}
\pgfplotsset{compat=1.18}

% Units, tables, floats, references
\usepackage[locale=DE,space-before-unit=true,per-mode=symbol,separate-uncertainty=true]{siunitx}
\usepackage{booktabs,multirow,array,tabularx,longtable}
\usepackage{caption,subcaption}
\usepackage{placeins}
\usepackage{enumitem}
\usepackage{xparse}

% Boxes
\usepackage[most]{tcolorbox}
\tcbuselibrary{breakable,skins,theorems}

% Hyperlinks should be loaded late
\usepackage[breaklinks=true,colorlinks=true,linkcolor=blue,urlcolor=blue,citecolor=blue]{hyperref}
\usepackage[nameinlink,noabbrev]{cleveref}

% Header and footer
\clearpairofpagestyles
\ihead{\CourseShortTitle}
\ohead{\DocumentKind}
\cfoot{\pagemark}
\pagestyle{scrheadings}

% Common list spacing
\setlist[enumerate]{itemsep=0.4em,topsep=0.4em}
\setlist[itemize]{itemsep=0.25em,topsep=0.35em}

% Float placement defaults
\captionsetup{font=small,labelfont=bf}

%%%%%%%%%%%%%%%%%%%%%%%%%%%%%%%%%%%%%%%%%%%%%%%%%%%%%%%%%%%%%%%%%%%
% Embedded file: includes/university_macros.tex
%%%%%%%%%%%%%%%%%%%%%%%%%%%%%%%%%%%%%%%%%%%%%%%%%%%%%%%%%%%%%%%%%%%
%%%%%%%%%%%%%%%%%%%%%%%%%%%%%%%%%%%%%%%%%%%%%%%%%%%%%%%%%%%%%%%%%%%
% General university LaTeX template -- mathematics, math, science, units macros
%%%%%%%%%%%%%%%%%%%%%%%%%%%%%%%%%%%%%%%%%%%%%%%%%%%%%%%%%%%%%%%%%%%

% Number sets
\newcommand{\R}{\mathbb{R}}
\newcommand{\N}{\mathbb{N}}
\newcommand{\Z}{\mathbb{Z}}
\newcommand{\Q}{\mathbb{Q}}
\newcommand{\C}{\mathbb{C}}

% Brackets and norms
\newcommand{\set}[1]{\left\{ #1 \right\}}
\newcommand{\abs}[1]{\left| #1 \right|}
\newcommand{\norm}[1]{\left\lVert #1 \right\rVert}
\newcommand{\avg}[1]{\left\langle #1 \right\rangle}
\newcommand{\paren}[1]{\left( #1 \right)}
\newcommand{\bracket}[1]{\left[ #1 \right]}
\newcommand{\ol}[1]{\overline{#1}}

% Differentials and derivatives
\newcommand{\dd}{\,\mathrm{d}}
\newcommand{\dv}[2]{\frac{\dd #1}{\dd #2}}
\newcommand{\pdv}[2]{\frac{\partial #1}{\partial #2}}
\newcommand{\pddv}[2]{\frac{\partial^2 #1}{\partial #2^2}}
\newcommand{\evalat}[2]{\left. #1 \right|_{#2}}

% Vectors and operators
\newcommand{\vect}[1]{\bm{#1}}
\newcommand{\unitvect}[1]{\hat{\bm{#1}}}
\newcommand{\grad}{\vect{\nabla}}
\newcommand{\divergence}{\grad\!\cdot}
\newcommand{\curl}{\grad\!\times}

% Frequent constants and factors
\newcommand{\half}{\frac{1}{2}}
\newcommand{\third}{\frac{1}{3}}
\newcommand{\twothirds}{\frac{2}{3}}
\newcommand{\hbarc}{\hbar c}
\newcommand{\alphaf}{\alpha}
\newcommand{\eV}{\si{eV}}
\newcommand{\MeV}{\si{MeV}}
\newcommand{\GeV}{\si{GeV}}
\newcommand{\TeV}{\si{TeV}}

% Nuclear and particle notation
\newcommand{\nuclide}[3]{^{#1}_{#2}\mathrm{#3}}
\newcommand{\isotope}[2]{^{#1}\mathrm{#2}}
\newcommand{\anti}[1]{\overline{#1}}
\newcommand{\pbar}{\anti{p}}
\newcommand{\nbar}{\anti{n}}
\newcommand{\ee}{e^-}
\newcommand{\ep}{e^+}
\newcommand{\mup}{\mu^+}
\newcommand{\mum}{\mu^-}
\newcommand{\nue}{\nu_e}
\newcommand{\numu}{\nu_\mu}
\newcommand{\nutau}{\nu_\tau}
\newcommand{\pionp}{\pi^+}
\newcommand{\pionm}{\pi^-}
\newcommand{\pionz}{\pi^0}
\newcommand{\kaonp}{K^+}
\newcommand{\kaonm}{K^-}
\newcommand{\kaonz}{K^0}

% Cross sections, luminosity, Mandelstam variables
\newcommand{\diffxs}{\frac{\dd\sigma}{\dd\Omega}}
\newcommand{\totalxs}{\sigma_{\mathrm{total}}}
\newcommand{\lum}{\mathcal{L}}
\newcommand{\smand}{s}
\newcommand{\tmand}{t}
\newcommand{\umand}{u}
\newcommand{\Ecm}{E_{\mathrm{CM}}}

% Common text shortcuts
\newcommand{\ie}{d.\,h.}
\newcommand{\eg}{z.\,B.}
\newcommand{\wrt}{bez\"uglich}

% Units not predefined by siunitx
\DeclareSIUnit{\barn}{b}
\DeclareSIUnit{\millibarn}{mb}
\DeclareSIUnit{\microbarn}{\micro b}
\DeclareSIUnit{\nanobarn}{nb}
\DeclareSIUnit{\fermi}{fm}
\DeclareSIUnit{\year}{a}

% Local shortcuts for this sheet
\newcommand{\Dzero}{\Delta^0}
\newcommand{\pivec}{\vect{p}_{\pi}}
\newcommand{\pstar}{p^{\ast}}
\newcommand{\Estar}{E^{\ast}}
\newcommand{\betacm}{\beta_{\mathrm{CM}}}
\newcommand{\gammacm}{\gamma_{\mathrm{CM}}}
\newcommand{\fourp}[4]{\paren{#1,#2,#3,#4}}

%%%%%%%%%%%%%%%%%%%%%%%%%%%%%%%%%%%%%%%%%%%%%%%%%%%%%%%%%%%%%%%%%%%
% Embedded file: includes/university_boxes.tex
%%%%%%%%%%%%%%%%%%%%%%%%%%%%%%%%%%%%%%%%%%%%%%%%%%%%%%%%%%%%%%%%%%%
%%%%%%%%%%%%%%%%%%%%%%%%%%%%%%%%%%%%%%%%%%%%%%%%%%%%%%%%%%%%%%%%%%%
% General university LaTeX template -- reusable tcolorbox environments
%%%%%%%%%%%%%%%%%%%%%%%%%%%%%%%%%%%%%%%%%%%%%%%%%%%%%%%%%%%%%%%%%%%

% Base style shared by most boxes.
\tcbset{
    unibox/.style={
        enhanced,
        breakable,
        sharp corners,
        boxrule=0.6pt,
        left=1.2mm,
        right=1.2mm,
        top=1mm,
        bottom=1mm,
        before skip=0.8em,
        after skip=0.8em,
        fonttitle=\bfseries
    }
}

% Calculation boxes may be nested. Optional argument overwrites title/options.
\newtcolorbox{calculationbox}[1][]{
    unibox,
    title=Rechnung,
    colback=black!2,
    colframe=black!55,
    #1
}

\newtcolorbox{sidecalcbox}[1][]{
    unibox,
    title=Nebenrechnung,
    colback=black!1,
    colframe=black!40,
    #1
}

\newtcolorbox{solutionbox}[1][]{
    unibox,
    title=L\"osung,
    colback=blue!2,
    colframe=blue!45!black,
    #1
}

\newtcolorbox{answerbox}[1][]{
    unibox,
    title=Antwort,
    colback=green!3,
    colframe=green!45!black,
    #1
}

\newtcolorbox{notebox}[1][]{
    unibox,
    title=Notiz,
    colback=yellow!6,
    colframe=yellow!45!black,
    #1
}

\newtcolorbox{definitionbox}[1][]{
    unibox,
    title=Definition,
    colback=cyan!3,
    colframe=cyan!45!black,
    #1
}

\newtcolorbox{warningbox}[1][]{
    unibox,
    title=Achtung,
    colback=red!3,
    colframe=red!55!black,
    #1
}

\newtcolorbox{summarybox}[1][]{
    unibox,
    title=Zusammenfassung,
    colback=violet!3,
    colframe=violet!50!black,
    #1
}

% Compact formula box for central formulas.
\newtcolorbox{formulabox}[1][]{
    unibox,
    title=Formel,
    colback=orange!4,
    colframe=orange!65!black,
    #1
}

% A task wrapper for homework sheets.
\newtcolorbox{taskbox}[1][]{
    unibox,
    title=Aufgabe,
    colback=white,
    colframe=black!75,
    #1
}

\newtcolorbox{presentationbox}[1][]{
    unibox,
    title=Tafelvortrag,
    colback=teal!3,
    colframe=teal!50!black,
    #1
}

%%%%%%%%%%%%%%%%%%%%%%%%%%%%%%%%%%%%%%%%%%%%%%%%%%%%%%%%%%%%%%%%%%%
% Embedded file: includes/university_theorems.tex
%%%%%%%%%%%%%%%%%%%%%%%%%%%%%%%%%%%%%%%%%%%%%%%%%%%%%%%%%%%%%%%%%%%
%%%%%%%%%%%%%%%%%%%%%%%%%%%%%%%%%%%%%%%%%%%%%%%%%%%%%%%%%%%%%%%%%%%
% General university LaTeX template -- theorem-like environments
%%%%%%%%%%%%%%%%%%%%%%%%%%%%%%%%%%%%%%%%%%%%%%%%%%%%%%%%%%%%%%%%%%%

\theoremstyle{definition}
\newtheorem{definition}{Definition}[section]
\newtheorem{example}{Beispiel}[section]

\theoremstyle{plain}
\newtheorem{satz}{Satz}[section]
\newtheorem{lemma}{Lemma}[section]

\theoremstyle{remark}
\newtheorem*{remark}{Bemerkung}

%%%%%%%%%%%%%%%%%%%%%%%%%%%%%%%%%%%%%%%%%%%%%%%%%%%%%%%%%%%%%%%%%%%
\begin{document}

\makeuniversitytitle
\tableofcontents
\newpage

%%%%%%%%%%%%%%%%%%%%%%%%%%%%%%%%%%%%%%%%%%%%%%%%%%%%%%%%%%%%%%%%%%%
% Homework content starts here
%%%%%%%%%%%%%%%%%%%%%%%%%%%%%%%%%%%%%%%%%%%%%%%%%%%%%%%%%%%%%%%%%%%

\section*{Vorbereitung f\"ur den Tafelvortrag}
\addcontentsline{toc}{section}{Vorbereitung f\"ur den Tafelvortrag}

\begin{summarybox}[title={Was man aus diesem Blatt lernen soll}]
Dieses Aufgabenblatt trainiert relativistische Kinematik in der Teilchenphysik. Die zentralen Werkzeuge sind
\begin{itemize}
    \item die invariante Masse beziehungsweise Mandelstam-Gr\"o\ss e $s$,
    \item die Umrechnung zwischen Laborsystem und Schwerpunktsystem,
    \item die Lorentztransformation von Energie und Impuls,
    \item die feste Energie- und Impulsaufteilung bei einem Zweik\"orperzerfall.
\end{itemize}
Wir verwenden nat\"urliche Einheiten $c=1$. Energien, Massen und Impulse werden deshalb in $\MeV$ oder $\GeV$ geschrieben. Wenn alle Komponenten eines Viererimpulses angegeben werden, ist die Konvention
\[
    p^\mu = (E,p_x,p_y,p_z)
\]
gemeint. Streng mit $c$ w\"are dies $p^\mu=(E/c,p_x,p_y,p_z)$; in Einheiten $c=1$ fallen beide Schreibweisen zusammen.
\end{summarybox}

\begin{presentationbox}[title={Roter Faden f\"ur den Vortrag}]
\begin{enumerate}
    \item Zuerst erkl\"aren: Nicht Energie oder Impuls einzeln sind allgemein invariant, sondern $s=(p_1+p_2)^2$.
    \item Bei ruhendem Target ausnutzen: $s=m_1^2+m_2^2+2m_2E_1$.
    \item Schwerpunktsystem-Geschwindigkeit aus Gesamtimpuls durch Gesamtenergie bestimmen: $\beta_{\mathrm{CM}}=p_{\mathrm{ges}}/E_{\mathrm{ges}}$.
    \item Beim Zerfall immer zuerst im Ruhesystem des Mutterteilchens rechnen. Danach mit Lorentztransformation ins Labor boosten.
    \item Bei Aufgabe 2 klar sagen: Der Boost ist entlang $+x$, die Zerfallsimpulse im Schwerpunktsystem liegen entlang $\pm y$. Daher bleibt $p_y$ unver\"andert und nur $E$ und $p_x$ mischen sich.
\end{enumerate}
\end{presentationbox}

\newpage

\section{Aufgabe 1: Alles ist relativ (egal)}

\begin{taskbox}
Betrachtet wird die Reaktion
\[
    \pionm + p \longrightarrow \Dzero \longrightarrow \pionz + n.
\]
Das Proton ruht im Laborsystem. Gesucht sind die Schwellenergie des einlaufenden $\pionm$, die Geschwindigkeit des Schwerpunktsystems bei dieser Schwellenergie und die maximal beziehungsweise minimal m\"ogliche Neutronenergie im Laborsystem nach dem Zerfall der $\Dzero$-Resonanz.
\end{taskbox}

\subsection{Teil (a): Schwellenergie f\"ur die Erzeugung der $\Dzero$-Resonanz}

\begin{presentationbox}
An der Tafel zuerst die Idee formulieren: Eine Resonanz kann erzeugt werden, wenn im Schwerpunktsystem mindestens ihre Ruheenergie vorhanden ist. Das hei\ss t hier
\[
    \sqrt{s}=m_\Delta.
\]
Dann $s$ im Laborsystem ausrechnen, weil dort das Proton ruht.
\end{presentationbox}

\begin{calculationbox}[title=Invariante Masse und Schwerpunktsenergie]
Die invariante Gr\"o\ss e ist das Quadrat des Gesamtviererimpulses:
\[
    s=(p_\pi+p_p)^2.
\]
Im Schwerpunktsystem verschwindet der Gesamtimpuls. Deshalb gilt dort
\[
    \sqrt{s}=E_{\mathrm{CM}}.
\]
An der Schwelle soll gerade die $\Dzero$-Resonanz erzeugt werden. Also ist
\[
    E_{\mathrm{CM}}=m_\Delta.
\]
\end{calculationbox}

\begin{calculationbox}[title=Auswertung im Laborsystem]
Das Proton ruht im Laborsystem. Also gilt
\[
    p_p=(m_p,0),
    \qquad
    p_\pi=(E_\pi,\vect{p}_\pi).
\]
Damit folgt
\begin{align*}
    s
    &= (p_\pi+p_p)^2 \\
    &= p_\pi^2+p_p^2+2p_\pi\cdot p_p \\
    &= m_\pi^2+m_p^2+2m_pE_\pi.
\end{align*}
An der Schwelle setzen wir $s=m_\Delta^2$:
\[
    m_\Delta^2=m_\pi^2+m_p^2+2m_pE_{\pi,\mathrm{thr}}.
\]
Aufl\"osen nach der Gesamtenergie des einlaufenden Pions liefert
\[
    E_{\pi,\mathrm{thr}}
    =\frac{m_\Delta^2-m_p^2-m_\pi^2}{2m_p}.
\]
Mit
\[
    m_\Delta=\SI{1232}{MeV},\qquad
    m_p=\SI{938,3}{MeV},\qquad
    m_{\pi^-}=\SI{139,6}{MeV}
\]
ergibt sich
\begin{align*}
    E_{\pi,\mathrm{thr}}
    &=\frac{(1232)^2-(938,3)^2-(139,6)^2}{2\cdot 938,3}\,\MeV \\
    &=\SI{329,28}{MeV}.
\end{align*}
Die kinetische Schwellenergie w\"are
\begin{align*}
    T_{\pi,\mathrm{thr}}
    &=E_{\pi,\mathrm{thr}}-m_\pi \\
    &=\SI{329,28}{MeV}-\SI{139,60}{MeV} \\
    &=\SI{189,68}{MeV}.
\end{align*}
\end{calculationbox}

\begin{answerbox}
Die Schwellenergie als Gesamtenergie des einlaufenden Pions ist
\[
    \boxed{E_{\pi,\mathrm{thr}}=\SI{329,28}{MeV}}.
\]
Die zugeh\"orige kinetische Schwellenergie ist
\[
    \boxed{T_{\pi,\mathrm{thr}}=\SI{189,68}{MeV}}.
\]
F\"ur die weiteren Rechnungen wird die Gesamtenergie $E_{\pi,\mathrm{thr}}$ verwendet.
\end{answerbox}

\subsection{Teil (b): Geschwindigkeit des Schwerpunktsystems relativ zum Labor}

\begin{presentationbox}
Die Geschwindigkeit des Schwerpunktsystems ist genau die Geschwindigkeit, mit der man in ein System boostet, in dem der Gesamtimpuls verschwindet. Im Labor ist der Gesamtimpuls nur der Pionimpuls, die Gesamtenergie ist $E_\pi+m_p$.
\end{presentationbox}

\begin{calculationbox}
F\"ur das System aus einlaufendem Pion und ruhendem Proton gilt
\[
    \betacm
    =\frac{\abs{\vect{p}_{\mathrm{ges}}}}{E_{\mathrm{ges}}}
    =\frac{p_\pi}{E_\pi+m_p}.
\]
Der Pionimpuls bei Schwellenergie folgt aus
\[
    E_\pi^2=p_\pi^2+m_\pi^2.
\]
Also ist
\begin{align*}
    p_{\pi,\mathrm{thr}}
    &=\sqrt{E_{\pi,\mathrm{thr}}^2-m_\pi^2} \\
    &=\sqrt{(329,28)^2-(139,6)^2}\,\MeV \\
    &=\SI{298,22}{MeV}.
\end{align*}
Damit folgt
\begin{align*}
    \betacm
    &=\frac{298,22}{329,28+938,3} \\
    &=0,23527.
\end{align*}
\end{calculationbox}

\begin{answerbox}
Das Schwerpunktsystem bewegt sich relativ zum Laborsystem mit
\[
    \boxed{\betacm=0,23527}
\]
beziehungsweise gerundet auf zwei Nachkommastellen mit
\[
    \boxed{v_{\mathrm{CM}}\approx 0,24c}.
\]
\end{answerbox}

\subsection{Teil (c): Maximale und minimale Neutronenergie im Laborsystem}

\begin{presentationbox}
Hier ist die entscheidende physikalische Idee: Im Schwerpunktsystem hat das Neutron immer dieselbe Energie $E_n^\ast$ und denselben Impulsbetrag $p_n^\ast$. Die Laborenergie h\"angt nur davon ab, ob das Neutron in Boostrichtung oder entgegen der Boostrichtung emittiert wird.
\end{presentationbox}

\begin{calculationbox}
Die Lorentztransformation f\"ur die Energie bei einem Boost entlang der Bewegungsrichtung des Schwerpunktsystems lautet
\[
    E_{\mathrm{lab}}
    =\gammacm\paren{E^\ast+\betacm p^\ast\cos\theta^\ast}.
\]
Dabei ist $\theta^\ast$ der Winkel zwischen Neutronimpuls im Schwerpunktsystem und Boostrichtung. Weiter gilt
\[
    \gammacm=\frac{1}{\sqrt{1-\betacm^2}}
    =\frac{1}{\sqrt{1-0,23527^2}}
    =1,02888.
\]
Gegeben sind
\[
    E_n^\ast=\SI{966,88}{MeV},\qquad
    p_n^\ast=\SI{228,19}{MeV}.
\]
Die maximale Energie erh\"alt man f\"ur $\cos\theta^\ast=+1$:
\begin{align*}
    E_{\max}
    &=\gammacm\paren{E_n^\ast+\betacm p_n^\ast} \\
    &=1,02888\paren{966,88+0,23527\cdot 228,19}\,\MeV \\
    &=\SI{1050,04}{MeV}.
\end{align*}
Die minimale Energie erh\"alt man f\"ur $\cos\theta^\ast=-1$:
\begin{align*}
    E_{\min}
    &=\gammacm\paren{E_n^\ast-\betacm p_n^\ast} \\
    &=1,02888\paren{966,88-0,23527\cdot 228,19}\,\MeV \\
    &=\SI{939,57}{MeV}.
\end{align*}
\end{calculationbox}

\begin{answerbox}
Die maximale Neutronenergie im Laborsystem ist
\[
    \boxed{E_{\max}=\SI{1050,04}{MeV}}.
\]
Die minimale Neutronenergie im Laborsystem ist
\[
    \boxed{E_{\min}=\SI{939,57}{MeV}}.
\]
Vorw\"arts emittierte Neutronen werden durch den Boost energetisch angehoben; r\"uckw\"arts emittierte Neutronen werden energetisch abgesenkt.
\end{answerbox}

\FloatBarrier
\newpage

\section{Aufgabe 2: Delta und invariante Masse}

\begin{taskbox}
Ein $\Dzero$-Baryon mit
\[
    m_{\Dzero}=\SI{1232}{MeV}/c^2
\]
und Gesamtenergie
\[
    E_\Delta=\SI{1,35}{GeV}
\]
fliegt im Laborsystem entlang der positiven $x$-Achse und zerf\"allt gem\"a\ss
\[
    \Dzero \longrightarrow p+\pionm.
\]
Im Schwerpunktsystem, also im Ruhesystem des $\Dzero$, sollen Proton und Pion entlang der $y$-Achse auseinanderfliegen.
\end{taskbox}

\subsection{Teil (a): Skizze beider Systeme}

\begin{calculationbox}[title=Skizze und Konvention]
Im Laborsystem bewegt sich das $\Dzero$ entlang $+x$. Im Schwerpunktsystem ruht das $\Dzero$ vor dem Zerfall. Wir w\"ahlen f\"ur die Rechnung
\[
    p \text{ in } +y\text{-Richtung},
    \qquad
    \pionm \text{ in } -y\text{-Richtung}.
\]

\begin{center}
\begin{tikzpicture}[>=Latex,scale=1.05]
    % Lab frame
    \node at (0,2.4) {Laborsystem};
    \draw[->] (-2.1,0) -- (2.2,0) node[right] {$x$};
    \draw[->] (0,-1.45) -- (0,1.65) node[above] {$y$};
    \draw[->,very thick] (-1.65,0) -- (-0.35,0) node[midway,above] {$\Dzero$};
    \fill (0,0) circle (1.5pt);
    \draw[->,thick] (0,0) -- (1.35,1.0) node[right] {$p$};
    \draw[->,thick] (0,0) -- (1.35,-1.0) node[right] {$\pionm$};

    % CM frame
    \begin{scope}[xshift=7cm]
        \node at (0,2.4) {Schwerpunktsystem};
        \draw[->] (-2.1,0) -- (2.2,0) node[right] {$x'$};
        \draw[->] (0,-1.45) -- (0,1.65) node[above] {$y'$};
        \fill (0,0) circle (1.5pt) node[below right] {$\Dzero$ ruht};
        \draw[->,thick] (0,0) -- (0,1.25) node[above] {$p$};
        \draw[->,thick] (0,0) -- (0,-1.25) node[below] {$\pionm$};
    \end{scope}
\end{tikzpicture}
\end{center}
\end{calculationbox}

\begin{answerbox}
Im Laborsystem hat das Gesamtsystem einen Impuls in $+x$-Richtung. Im Schwerpunktsystem ruht das $\Dzero$, und die Zerfallsprodukte haben entgegengesetzte Impulse entlang der $y$-Achse.
\end{answerbox}

\subsection{Teil (b): Geschwindigkeit des $\Dzero$ im Laborsystem}

\begin{calculationbox}
F\"ur ein einzelnes relativistisches Teilchen gilt
\[
    E^2=p^2+m^2,
    \qquad
    \beta=\frac{p}{E}.
\]
Mit
\[
    E_\Delta=\SI{1,35}{GeV},
    \qquad
    m_\Delta=\SI{1,232}{GeV}
\]
ergibt sich zuerst
\begin{align*}
    p_\Delta
    &=\sqrt{E_\Delta^2-m_\Delta^2} \\
    &=\sqrt{(1,35)^2-(1,232)^2}\,\GeV \\
    &=\SI{0,55198}{GeV}.
\end{align*}
Damit folgt
\begin{align*}
    \beta
    &=\frac{p_\Delta}{E_\Delta} \\
    &=\frac{0,55198}{1,35} \\
    &=0,40887.
\end{align*}
Der Lorentzfaktor ist
\[
    \gamma=\frac{E_\Delta}{m_\Delta}
    =\frac{1,35}{1,232}
    =1,09578.
\]
\end{calculationbox}

\begin{answerbox}
Das $\Dzero$ bewegt sich im Laborsystem mit
\[
    \boxed{\beta=0,40887}
\]
beziehungsweise gerundet
\[
    \boxed{v\approx 0,41c}.
\]
\end{answerbox}

\subsection{Teil (c): Viererimpulse im Schwerpunktsystem}

\begin{presentationbox}
Beim Zweik\"orperzerfall im Ruhesystem des Mutterteilchens sind die Energien der Tochterteilchen kinematisch fest. Der Impulsbetrag ist f\"ur beide gleich gro\ss, aber die Richtungen sind entgegengesetzt.
\end{presentationbox}

\begin{formulabox}[title=Zweik\"orperzerfall]
F\"ur $M\to 1+2$ im Ruhesystem von $M$ gilt
\[
    E_1^\ast=\frac{M^2+m_1^2-m_2^2}{2M},
    \qquad
    E_2^\ast=\frac{M^2+m_2^2-m_1^2}{2M}.
\]
\end{formulabox}

\begin{calculationbox}
Wir verwenden
\[
    m_\Delta=\SI{1,232}{GeV},
    \qquad
    m_p=\SI{0,9383}{GeV},
    \qquad
    m_\pi=\SI{0,1396}{GeV}.
\]
Damit gilt
\begin{align*}
    E_p^\ast
    &=\frac{m_\Delta^2+m_p^2-m_\pi^2}{2m_\Delta} \\
    &=\frac{1,232^2+0,9383^2-0,1396^2}{2\cdot 1,232}\,\GeV \\
    &=\SI{0,96540}{GeV},\\[0.5em]
    E_\pi^\ast
    &=\frac{m_\Delta^2+m_\pi^2-m_p^2}{2m_\Delta} \\
    &=\frac{1,232^2+0,1396^2-0,9383^2}{2\cdot 1,232}\,\GeV \\
    &=\SI{0,26660}{GeV}.
\end{align*}
Der gemeinsame Impulsbetrag ist
\begin{align*}
    p^\ast
    &=\sqrt{(E_p^\ast)^2-m_p^2} \\
    &=\sqrt{(0,96540)^2-(0,9383)^2}\,\GeV \\
    &=\SI{0,22713}{GeV}.
\end{align*}
Da wir $p$ in $+y$-Richtung und $\pi^-$ in $-y$-Richtung w\"ahlen, sind die Viererimpulse im Schwerpunktsystem
\[
    p_p^{\ast\mu}=\fourp{0,96540}{0}{+0,22713}{0}\,\GeV,
\]
\[
    p_\pi^{\ast\mu}=\fourp{0,26660}{0}{-0,22713}{0}\,\GeV.
\]
\end{calculationbox}

\begin{answerbox}
Die Viererimpulse im Schwerpunktsystem sind
\[
    \boxed{p_p^{\ast\mu}=\fourp{0,96540}{0}{+0,22713}{0}\,\GeV}
\]
und
\[
    \boxed{p_\pi^{\ast\mu}=\fourp{0,26660}{0}{-0,22713}{0}\,\GeV}.
\]
Die Summe ist
\[
    p_p^{\ast\mu}+p_\pi^{\ast\mu}
    =\fourp{1,23200}{0}{0}{0}\,\GeV,
\]
also genau der Viererimpuls des ruhenden $\Dzero$.
\end{answerbox}

\subsection{Teil (d): Viererimpulse im Laborsystem}

\begin{presentationbox}
Der Boost geht vom Ruhesystem des $\Dzero$ ins Laborsystem. Weil das $\Dzero$ im Labor in $+x$-Richtung fliegt, steht in der Boostmatrix das Pluszeichen bei $\gamma\beta$.
\end{presentationbox}

\begin{calculationbox}[title=Lorentztransformation]
F\"ur den Boost aus dem Schwerpunktsystem ins Laborsystem gilt
\[
\begin{pmatrix}
E\\ p_x\\ p_y\\ p_z
\end{pmatrix}_{\!\mathrm{lab}}
=
\begin{pmatrix}
\gamma & \gamma\beta & 0 & 0\\
\gamma\beta & \gamma & 0 & 0\\
0&0&1&0\\
0&0&0&1
\end{pmatrix}
\begin{pmatrix}
E^\ast\\ p_x^\ast\\ p_y^\ast\\ p_z^\ast
\end{pmatrix}.
\]
Da $p_x^\ast=0$ gilt, vereinfacht sich das zu
\[
    E_{\mathrm{lab}}=\gamma E^\ast,
    \qquad
    p_{x,\mathrm{lab}}=\gamma\beta E^\ast,
    \qquad
    p_{y,\mathrm{lab}}=p_y^\ast.
\]
Mit
\[
    \beta=0,40887,
    \qquad
    \gamma=1,09578
\]
folgt f\"ur das Proton
\begin{align*}
    E_p&=1,09578\cdot 0,96540=\SI{1,05786}{GeV},\\
    p_{p,x}&=1,09578\cdot 0,40887\cdot 0,96540=\SI{0,43253}{GeV},\\
    p_{p,y}&=+\SI{0,22713}{GeV}.
\end{align*}
Also
\[
    p_p^\mu=\fourp{1,05786}{0,43253}{+0,22713}{0}\,\GeV.
\]
F\"ur das Pion folgt analog
\begin{align*}
    E_\pi&=1,09578\cdot 0,26660=\SI{0,29214}{GeV},\\
    p_{\pi,x}&=1,09578\cdot 0,40887\cdot 0,26660=\SI{0,11945}{GeV},\\
    p_{\pi,y}&=-\SI{0,22713}{GeV}.
\end{align*}
Also
\[
    p_\pi^\mu=\fourp{0,29214}{0,11945}{-0,22713}{0}\,\GeV.
\]
\end{calculationbox}

\begin{calculationbox}[title=Kontrolle durch Summenbildung]
Die Summe der beiden Viererimpulse ist
\begin{align*}
    p_p^\mu+p_\pi^\mu
    &=\fourp{1,05786+0,29214}{0,43253+0,11945}{0,22713-0,22713}{0}\,\GeV\\
    &=\fourp{1,35000}{0,55198}{0}{0}\,\GeV.
\end{align*}
Das stimmt mit dem Viererimpuls des $\Dzero$ im Laborsystem \"uberein:
\[
    p_\Delta^\mu=(E_\Delta,p_\Delta,0,0)
    =\fourp{1,35000}{0,55198}{0}{0}\,\GeV.
\]
\end{calculationbox}

\begin{answerbox}
Die Viererimpulse im Laborsystem sind
\[
    \boxed{p_p^\mu=\fourp{1,05786}{0,43253}{+0,22713}{0}\,\GeV}
\]
und
\[
    \boxed{p_\pi^\mu=\fourp{0,29214}{0,11945}{-0,22713}{0}\,\GeV}.
\]
Die Kontrolle ergibt
\[
    \boxed{p_p^\mu+p_\pi^\mu=\fourp{1,35000}{0,55198}{0}{0}\,\GeV}.
\]
\end{answerbox}

\subsection{Teil (e): Winkel und \"Offnungswinkel im Laborsystem}

\begin{calculationbox}
Der Winkel zur positiven $x$-Achse ergibt sich aus den Impulskomponenten:
\[
    \tan\alpha=\frac{\abs{p_y}}{p_x}.
\]
F\"ur das Proton folgt
\begin{align*}
    \alpha_p
    &=\arctan\paren{\frac{0,22713}{0,43253}} \\
    &=27,70^\circ.
\end{align*}
F\"ur das Pion folgt
\begin{align*}
    \alpha_\pi
    &=\arctan\paren{\frac{0,22713}{0,11945}} \\
    &=62,26^\circ.
\end{align*}
Da Proton und Pion auf verschiedenen Seiten der $x$-Achse fliegen, ist der \"Offnungswinkel die Summe:
\begin{align*}
    \theta_{\mathrm{open}}
    &=\alpha_p+\alpha_\pi \\
    &=27,70^\circ+62,26^\circ \\
    &=89,97^\circ.
\end{align*}
\end{calculationbox}

\begin{answerbox}
Der Winkel des Protons zur $x$-Achse ist
\[
    \boxed{\alpha_p=27,70^\circ}.
\]
Der entsprechende Winkel des Pions zur $x$-Achse ist
\[
    \boxed{\alpha_\pi=62,26^\circ}.
\]
Der \"Offnungswinkel zwischen beiden Bahnen ist
\[
    \boxed{\theta_{\mathrm{open}}=89,97^\circ}.
\]
\end{answerbox}

\subsection{Teil (f): Vergleich der Geschwindigkeiten von Proton und Pion}

\begin{presentationbox}
Qualitativ kann man zuerst ohne Rechnen argumentieren: Im Schwerpunktsystem haben beide denselben Impulsbetrag. Das leichtere Teilchen hat bei gleichem Impuls die gr\"o\ss ere Geschwindigkeit. Das ist das Pion. Im Laborsystem wird zwar in $x$-Richtung geboostet, aber das Pion bleibt deutlich schneller.
\end{presentationbox}

\begin{calculationbox}[title=Qualitative Begr\"undung]
Im Schwerpunktsystem haben Proton und Pion betragsm\"a\ss ig denselben Impuls, weil sie aus einem ruhenden Mutterteilchen in entgegengesetzte Richtungen zerfallen. F\"ur ein relativistisches Teilchen gilt
\[
    \frac{v}{c}=\frac{p}{E}.
\]
Bei gleichem Impuls hat das leichtere Teilchen die kleinere Energie und damit das gr\"o\ss ere $p/E$. Da
\[
    m_\pi \ll m_p,
\]
ist das Pion im Schwerpunktsystem schneller.

Auch im Laborsystem bleibt das Pion schneller. Der Boost in $+x$-Richtung ver\"andert die Impulskomponenten, aber das Pion hat weiterhin den deutlich gr\"o\ss eren Geschwindigkeitsbetrag.
\end{calculationbox}

\begin{sidecalcbox}[title=Kontrolle durch direkte Rechnung]
Mit $v/c=\abs{\vect{p}}/E$ erh\"alt man im Laborsystem
\begin{align*}
    \frac{v_p}{c}
    &=\frac{\sqrt{0,43253^2+0,22713^2}}{1,05786}
    =0,46182,\\[0.5em]
    \frac{v_\pi}{c}
    &=\frac{\sqrt{0,11945^2+0,22713^2}}{0,29214}
    =0,87844.
\end{align*}
\end{sidecalcbox}

\begin{answerbox}
Sowohl im Schwerpunktsystem als auch im Laborsystem ist das Pion schneller:
\[
    \boxed{v_\pi>v_p}.
\]
Im Laborsystem gilt zur Kontrolle n\"aherungsweise
\[
    \boxed{v_p\approx 0,46c,\qquad v_\pi\approx 0,88c}.
\]
\end{answerbox}

\FloatBarrier
\newpage

\section*{Lernzusammenfassung}
\addcontentsline{toc}{section}{Lernzusammenfassung}

\begin{summarybox}[title={Formeln, die man sicher beherrschen sollte}]
\begin{align*}
    s&=(p_1+p_2)^2,\\
    s_{\text{ruhendes Target}}&=m_1^2+m_2^2+2m_2E_1,\\
    E_{\mathrm{thr}}&=\frac{M^2-m_2^2-m_1^2}{2m_2},\\
    \beta_{\mathrm{CM}}&=\frac{p_{\mathrm{ges}}}{E_{\mathrm{ges}}},\\
    E_{\mathrm{lab}}&=\gamma(E^\ast+\beta p_x^\ast),\\
    p_{x,\mathrm{lab}}&=\gamma(p_x^\ast+\beta E^\ast),\\
    p_y&=p_y^\ast,\qquad p_z=p_z^\ast \quad \text{bei Boost entlang } x.
\end{align*}
\end{summarybox}

\begin{warningbox}[title={Typische Fehlerquellen}]
\begin{itemize}
    \item Schwellenergie kann Gesamtenergie oder kinetische Energie bedeuten. In der invarianten Rechnung ist zun\"achst $E_{\pi,\mathrm{thr}}$ die Gesamtenergie.
    \item Beim Lorentzboost muss die Richtung stimmen. Vom Ruhesystem ins Labor bei Bewegung in $+x$ verwendet man $+\gamma\beta$.
    \item Ein Boost entlang $x$ ver\"andert $E$ und $p_x$, aber nicht $p_y$ und $p_z$.
    \item Beim \"Offnungswinkel in Aufgabe 2 liegen Proton und Pion auf verschiedenen Seiten der $x$-Achse. Deshalb werden die beiden Einzelwinkel addiert.
\end{itemize}
\end{warningbox}

\begin{summarybox}[title={Endergebnisse}]
\begin{itemize}
    \item Aufgabe 1(a): $E_{\pi,\mathrm{thr}}=\SI{329,28}{MeV}$, $T_{\pi,\mathrm{thr}}=\SI{189,68}{MeV}$.
    \item Aufgabe 1(b): $\beta_{\mathrm{CM}}=0,23527\approx 0,24$.
    \item Aufgabe 1(c): $E_{\max}=\SI{1050,04}{MeV}$, $E_{\min}=\SI{939,57}{MeV}$.
    \item Aufgabe 2(b): $\beta_\Delta=0,40887\approx 0,41$, $\gamma=1,09578$.
    \item Aufgabe 2(c): $p_p^{\ast\mu}=\fourp{0,96540}{0}{+0,22713}{0}\,\GeV$, $p_\pi^{\ast\mu}=\fourp{0,26660}{0}{-0,22713}{0}\,\GeV$.
    \item Aufgabe 2(d): $p_p^\mu=\fourp{1,05786}{0,43253}{+0,22713}{0}\,\GeV$, $p_\pi^\mu=\fourp{0,29214}{0,11945}{-0,22713}{0}\,\GeV$.
    \item Aufgabe 2(e): $\alpha_p=27,70^\circ$, $\alpha_\pi=62,26^\circ$, $\theta_{\mathrm{open}}=89,97^\circ$.
    \item Aufgabe 2(f): $v_\pi>v_p$, im Labor ungef\"ahr $v_p=0,46c$ und $v_\pi=0,88c$.
\end{itemize}
\end{summarybox}

\end{document}
